\documentclass[11pt,a4paper]{article}
\usepackage[T1]{fontenc}
\usepackage[utf8]{inputenc}
\usepackage{amsmath,amssymb,xcolor}
\usepackage[margin=25mm]{geometry}
\usepackage[colorlinks=true,urlcolor=blue]{hyperref}
\definecolor{nexorablue}{HTML}{2459B5}
\setlength{\parindent}{0pt}
\setlength{\parskip}{7pt}
\setlength{\emergencystretch}{2em}
\renewcommand{\familydefault}{\sfdefault}
\begin{document}

{\small\bfseries\color{nexorablue} NEXORA CPGE}\par
{\small MPSI / NEX-CIR-06 / Corrigé}\par
{\LARGE\bfseries Deux nœuds et un bilan complet}\par
\vspace{0.25em}\hrule\vspace{0.6em}

\section*{Équations de nœuds}
En $C$, $V_B-V_C=V_C/2$, soit $V_C=2V_B/3$. En $B$, $(12-V_B)/2=V_B/2+V_B-V_C$. On trouve $V_B=4{,}5\,\mathrm V$ et $V_C=3{,}0\,\mathrm V$.

\section*{Courants et puissances}
Dans l’ordre $AB,BC,B0,C0$, les courants valent $3{,}75$, $1{,}50$, $2{,}25$ et $1{,}50\,\mathrm{mA}$. Les puissances dissipées sont $28{,}125$, $2{,}25$, $10{,}125$ et $4{,}50\,\mathrm{mW}$, soit $45\,\mathrm{mW}$ au total. La source fournit $12\times3{,}75=45\,\mathrm{mW}$.


\end{document}
